A point path follows one chosen input through a sequence of motions.  The
projective operation $c/z$ admits no such motion of the value, and the
observation of Subsection~\ref{subsec:p0-point-paths} was that its matrix has a
nonzero lower-left entry.  There is nevertheless a geometry attached to it, and
it is obtained by reversing the direction of the reading: instead of pushing a
point forward, pull an entire output condition backward through the same maps.
The upper half-plane makes this dual reading elementary, because M\"obius maps
carry generalized circles to generalized circles.

We call the resulting picture \emph{ripple geometry}: the family of conditions
obtained by pulling output conditions back through a projective arithmetic
operator.  The name is descriptive, and
Subsection~\ref{subsec:p0-ripple-status} states precisely what it does and does
not assert.

\subsection{Poles and standard horocycles}

Let
\begin{equation}\label{eq:p0-mobius-map}
  F(z)=\frac{Az+B}{Cz+D}
\end{equation}
be a real fractional linear map with
\[
  \Delta=AD-BC>0.
\]
The positivity convention ensures that $F$ maps the upper half-plane to itself.
The projective pole is
\[
  p_F=F^{-1}(\infty)
  =\begin{cases}
     -D/C,&C\ne0,\\
     \infty,&C=0.
   \end{cases}
\]
For $t>0$, denote the standard horocycle based at infinity by
\[
  \mathcal H_t=\{w\in\mathbb H:\im w=t\}.
\]

\begin{definition}[Ripple pencil]
\label{def:p0-ripple-pencil}
The \emph{ripple pencil} of $F$ is the one-parameter family
\[
  \mathcal R(F)
  =\{F^{-1}(\mathcal H_t):t>0\}.
\]
It is a family of horocycles based at the pole $p_F$.  The word ``ripple'' is a
visual name for this standard M\"obius pullback family.
\end{definition}

\begin{theorem}[Line-to-circle transition]
\label{thm:p0-ripple-pencil-formula}
Let $F$ be as in \eqref{eq:p0-mobius-map}.

\begin{enumerate}
  \item If $C=0$, every member of $\mathcal R(F)$ is a horizontal line; the
  pencil is based at $\infty$.

  \item If $C\ne0$, the member with parameter $t$ is the Euclidean circle
  \begin{equation}\label{eq:p0-ripple-circle}
    \left(x+\frac DC\right)^2
    +\left(y-\frac{\Delta}{2tC^2}\right)^2
    =\left(\frac{\Delta}{2tC^2}\right)^2.
  \end{equation}
  All these circles are tangent to the real boundary at the common point
  $p_F=-D/C$ and are nested as $t$ varies.
\end{enumerate}
\end{theorem}

\begin{proof}
For $z=x+iy$, direct calculation gives the standard identity
\begin{equation}\label{eq:p0-imaginary-mobius}
  \im F(z)=\frac{\Delta y}{|Cz+D|^2}.
\end{equation}
The inverse image of $\mathcal H_t$ is therefore
\[
  \Delta y=t\bigl((Cx+D)^2+C^2y^2\bigr).
\]
If $C=0$, this reduces to a horizontal line.  If $C\ne0$, divide by
$tC^2$ and complete the square in $y$:
\[
  \left(x+\frac DC\right)^2+y^2
  -\frac{\Delta}{tC^2}y=0,
\]
which is exactly \eqref{eq:p0-ripple-circle}.  The center lies vertically
above $-D/C$ at a height equal to the radius, so the circle is tangent to the
boundary there.  The radius is proportional to $t^{-1}$, proving nesting.
\end{proof}

The theorem gives a geometric interpretation of the lower-left matrix entry.
The affine condition $C=0$ places the pole at infinity and degenerates the
ripple circles to parallel lines.  A projective step with $C\ne0$ moves the pole
to a finite boundary point and curves the same horocyclic family into visible
ripples.

\begin{example}[Arithmetic inversion]
\label{ex:p0-inversion-ripples}
For
\[
  J(z)=-\frac1z,
  \qquad
  J\leftrightarrow
  \begin{pmatrix}0&-1\\1&0\end{pmatrix},
\]
one has $p_J=0$ and
\begin{equation}\label{eq:p0-inversion-circles}
  x^2+\left(y-\frac1{2t}\right)^2
  =\left(\frac1{2t}\right)^2.
\end{equation}
Thus inversion turns the horizontal horocycles based at infinity into a nested
family of circles tangent at zero.  The map exchanges the two ideal points
$0$ and $\infty$.
\end{example}

\begin{figure}[t]
\centering
\begin{tikzpicture}[
  x=1.05cm,y=1.05cm,
  axis/.style={-{Latex[length=2mm]},thick,black!70},
  ripple/.style={thick,blue!60!black},
  geod/.style={dashed,gray!55},
  every node/.style={font=\small}
]
  \fill[blue!2] (-4.2,0) rectangle (4.2,4.0);
  \draw[axis] (-4.35,0)--(4.45,0) node[right] {$x$};
  \draw[axis] (0,-0.08)--(0,4.2) node[above] {$y$};
  \foreach \r in {0.45,0.75,1.15,1.7}
    {\draw[ripple] (0,\r) circle (\r);}
  \draw[geod] (0,0)--(0,3.75);
  \fill[orange!80!black] (0,0) circle (2.2pt)
    node[below right] {pole $p_J=0$};
  \node[blue!55!black,anchor=west] at (1.85,2.75)
    {$J^{-1}(\mathcal H_t)$};
  \node[anchor=west,font=\scriptsize] at (1.85,2.3)
    {larger circles correspond to smaller $t$};
\end{tikzpicture}
\caption{The ripple pencil of inversion.  Every circle is a horocycle tangent
at the projective pole $0$; the affine horizontal-line pencil is recovered
when the pole is moved to infinity.}
\label{fig:p0-inversion-ripple-pencil}
\end{figure}

\begin{remark}[Which upper half-plane, and with which metric]
\label{rem:p0-ripple-normalization}
The complex coordinate used throughout this section is the Möbius continuation
chart, and it is not literally the point coordinate of
Section~\ref{sec:p0-aes}.  With the metric $g_{\mu,\lambda}$ of
\eqref{eq:p0-e0-data}, the chart in which the horocycles are the Euclidean
horizontal lines and the metric is the standard one of curvature $-1$ is
\[
  Z=\frac{\lambda}{\mu}x+iy,
\]
as follows from \eqref{eq:scaled-standard-hyperbolic-metric}.  When
$\mu=\lambda$ this is the identity, which is why the two roles of $x+iy$ are
usually identified in the literature; here they are kept apart, and all
Möbius statements below are made in the $Z$-chart.  Only when a statement
concerns the assignment $a$ rather than the metric do we return to $(x,y)$.
\end{remark}

\subsection{Zero--pole pencils in the complex chart}

There is a second elementary pencil attached to the same operator.  Suppose
$A,C\ne0$ and write its zero and pole as
\[
  q_F=-\frac BA,
  \qquad
  p_F=-\frac DC.
\]
Then the output level set $|F(z)|=r$ is equivalent to
\begin{equation}\label{eq:p0-apollonius-pencil}
  |z-q_F|
  =r\left|\frac CA\right|\,|z-p_F|.
\end{equation}
For $r|C/A|\ne1$ these are Apollonius circles; for $r|C/A|=1$ the locus is the
perpendicular bisector line of the segment $q_Fp_F$, that is, a generalized
circle of infinite radius.  The horocyclic
ripple pencil and the Apollonius pencil encode different output observables,
$\im F$ and $|F|$, but both show that a projective one-hole operator naturally
carries more than a point orbit: it carries a zero, a pole, and a family of
generalized circles organized by them.

\subsection{Ripple histories of right-expanded expressions}

Let a right-expanded expression determine nondegenerate real M\"obius sections
$F_1,\ldots,F_n$ of \emph{positive} determinant, in chronological order, and
write
\[
  G_k=F_k\circ\cdots\circ F_1,
  \qquad
  G_k(z)=\frac{A_kz+B_k}{C_kz+D_k}.
\]
At stage $k$, the arithmetic value of a chosen input follows the point orbit
$G_k(z_0)$.  Simultaneously, the output horocycles pull back to
\[
  \mathcal R_k=\mathcal R(G_k).
\]
The tangency point
\[
  p_k=G_k^{-1}(\infty)
  =-\frac{D_k}{C_k}
\]
when $C_k\ne0$ is the current pole of the prefix operator.

Thus a general right-expansion history should not be pictured as one fixed
family of concentric circles.  Each prefix has its own ripple pencil, and the
common tangency point of that pencil may move with $k$.  The appropriate
geometric state is
\[
  \bigl(G_k,\,G_k(z_0),\,p_k,\,\mathcal R_k\bigr):
\]
the operator, the selected point, the pole, and the pulled-back wavefront
family.

\begin{proposition}[Covariant--contravariant duality]
\label{prop:p0-path-ripple-duality}
For every prefix operator $G_k$ and every output horocycle $\mathcal H_t$,
\[
  z_0\in G_k^{-1}(\mathcal H_t)
  \quad\Longleftrightarrow\quad
  G_k(z_0)\in\mathcal H_t.
\]
Hence the point path and the ripple history are two readings of the same
operator word: one pushes the input forward, the other pulls the output
condition backward.
\end{proposition}

\begin{proof}
This is the defining property of inverse image.
\end{proof}

The elementary nature of the proposition is precisely its value.  No new
geometry is inserted between the two pictures.  Their unity is already present
in the action of one M\"obius matrix.

\subsection{The status of ripple geometry}
\label{subsec:p0-ripple-status}

The preceding constructions are elementary and proved.  The name
``ripple geometry'' is not, and it is important to separate the two.

\begin{itemize}
  \item \emph{What is proved.}  A real orientation-preserving M\"obius map
  pulls the standard horocycle $\mathcal H_t$ back to a line (when its pole is
  infinity) or to a circle tangent to the boundary at its pole, and these
  circles are nested; the pole is the inverse image of infinity; the same
  prefix matrix determines the point orbit, the pole, the fixed points, and the
  ripple pencil.
  \item \emph{What is descriptive.}  ``Ripple geometry'' names the dual reading
  in which the operator acts on conditions instead of points.  It does not
  assert a new isomorphism class of hyperbolic or projective geometry, a new
  invariant of arithmetic words, or the existence of a projective arithmetic
  expression space.
  \item \emph{What is open.}  Whether the ripple family, taken together with
  its moving pole, is intrinsic to the arithmetic word rather than to a chosen
  upper-half-plane chart; which path or ripple data descend through relations
  among history words; and whether a punctured model in the sense of
  Section~\ref{sec:p0-holed} carries a ripple structure that survives at the
  puncture.
\end{itemize}

The reason for keeping the name is that the dual reading is what makes
division homogeneous with the other operations at the level of *conditions*:
addition, subtraction, and multiplication by a constant move the value and
hence move the contours of the assignment, while second-slot division moves the
contours without moving the value affinely.  The reason for keeping it
provisional is that no intrinsic formulation of that observation is available
yet; the current statement of it depends on the standard horocycle
$\mathcal H_t$ and therefore on one chart.
